
\documentclass{article}
\usepackage{colortbl}
\usepackage{makecell}
\usepackage{multirow}
\usepackage{supertabular}

\begin{document}

\newcounter{utterance}

\centering \large Interaction Transcript for game `taboo', experiment `low\_en', episode 0 with reflexion{-}llama{-}3.1{-}8b{-}instruct.
\vspace{24pt}

{ \footnotesize  \setcounter{utterance}{1}
\setlength{\tabcolsep}{0pt}
\begin{supertabular}{c@{$\;$}|p{.15\linewidth}@{}p{.15\linewidth}p{.15\linewidth}p{.15\linewidth}p{.15\linewidth}p{.15\linewidth}}
    \# & $\;$A & \multicolumn{4}{c}{Game Master} & $\;\:$B\\
    \hline

    \theutterance \stepcounter{utterance}  
    & & \multicolumn{4}{p{0.6\linewidth}}{
        \cellcolor[rgb]{0.9,0.9,0.9}{
            \makecell[{{p{\linewidth}}}]{
                \texttt{\tiny{[P1$\langle$GM]}}
                \texttt{You have attempted this same task before and did not succeed. These are the plans you wrote after those attempts.} \\
\texttt{Trial \#1: New plan: I failed to provide a description that would lead the other player to guess the target word within three attempts. In my last clue, I described the target word as "something that has just been created or introduced," which is a characteristic of the target word "novelty." However, I should have been more careful not to use related words or characteristics that are too specific to the target word. Next time, I will try to provide a more general description that captures the essence of the target word without giving away too much information. For example, I could have described it as "something that adds value or interest to an experience."} \\
\\ 
\texttt{Trial \#2: New plan: I failed to provide a description that would lead the other player to guess the target word within three attempts. In my last clue, I described the target word as "something that has just been created or introduced," which is a characteristic of the target word "novelty." However, I should have been more careful not to use related words or characteristics that are too specific to the target word. Next time, I will try to provide a more general description that captures the essence of the target word without giving away too much information. For example, I could have described it as "something that adds value or interest to an experience." } \\
\\ 
\texttt{To achieve this, I will focus on providing a description that highlights the benefits or effects of the target word, rather than its characteristics or related words. I will also try to use more abstract language to avoid giving away too much information. Additionally, I will pay closer attention to the related words provided and avoid using language that is too similar to them.} \\
\\ 
\texttt{Trial \#3: New plan: I failed to provide a description that would lead the other player to guess the target word within three attempts. In my last clue, I described the target word as "something that brings a new perspective to a familiar situation," which is a characteristic of the target word "novelty." However, I should have been more careful not to use related words or characteristics that are too specific to the target word. Next time, I will try to provide a more general description that captures the essence of the target word without giving away too much information. } \\
\\ 
\texttt{To achieve this, I will focus on providing a description that highlights the benefits or effects of the target word, rather than its characteristics or related words. I will also try to use more abstract language to avoid giving away too much information. Additionally, I will pay closer attention to the related words provided and avoid using language that is too similar to them. } \\
\\ 
\texttt{In particular, I will try to provide a description that is more focused on the emotional or experiential aspects of the target word, rather than its literal meaning. For example, I could have described it as "something that creates a sense of wonder or excitement." This approach will allow me to provide a more general description that captures the essence of the target word without giving away too much information.} \\
\\ 
\\ 
\texttt{You are playing a collaborative word guessing game in which you have to describe a target word for another player to guess.} \\
\\ 
\texttt{Rules:} \\
\texttt{(a) You have to reply in the form: CLUE: <some text>. Guesses from the other player will start with GUESS.} \\
\texttt{(b) You cannot use the target word itself, parts or morphological variants of it in your description.} \\
\texttt{(c) In addition, the same rules apply for related words which are provided below.} \\
\\ 
\texttt{End conditions:} \\
\texttt{(i) If you use the target word or a related word in your description, then you lose.} \\
\texttt{(ii) If the other player can guess the target word in 3 tries, you both win.} \\
\\ 
\texttt{Let us start.} \\
\\ 
\texttt{This is the target word that you need to describe and that the other player needs to guess:} \\
\\ 
\texttt{novelty} \\
\\ 
\texttt{Related words are:} \\
\\ 
\texttt{{-} oddity} \\
\texttt{{-} originality} \\
\texttt{{-} freshness} \\
\\ 
\texttt{Important: You are under time pressure, give short descriptions that are to the point!} \\
            }
        }
    }
    & \\ \\

    \theutterance \stepcounter{utterance}  
    & \multicolumn{4}{p{0.6\linewidth}}{
        \cellcolor[rgb]{0.9,0.9,0.9}{
            \makecell[{{p{\linewidth}}}]{
                \texttt{\tiny{[P1$\rangle$GM]}}
                \texttt{CLUE: Something that adds a unique twist to familiar things, making them more engaging and interesting.} \\
            }
        }
    }
    & & \\ \\

    \theutterance \stepcounter{utterance}  
    & & & \multicolumn{2}{p{0.3\linewidth}}{
        \cellcolor[rgb]{0.9,0.9,0.9}{
            \makecell[{{p{\linewidth}}}]{
                \texttt{\tiny{[GM$|$GM]}}
                \texttt{continue} \\
            }
        }
    }
    & & \\ \\

    \theutterance \stepcounter{utterance}  
    & & & \multicolumn{2}{p{0.3\linewidth}}{
        \cellcolor[rgb]{0.9,0.9,0.9}{
            \makecell[{{p{\linewidth}}}]{
                \texttt{\tiny{[GM$|$GM]}}
                \texttt{something that adds a unique twist to familiar things making them more engaging and interesting} \\
            }
        }
    }
    & & \\ \\

    \theutterance \stepcounter{utterance}  
    & & \multicolumn{4}{p{0.6\linewidth}}{
        \cellcolor[rgb]{0.9,0.9,0.9}{
            \makecell[{{p{\linewidth}}}]{
                \texttt{\tiny{[GM$\rangle$P2]}}
                \texttt{You have attempted this same task before and did not succeed. These are the plans you wrote after those attempts.} \\
\texttt{Trial \#1: New plan: I failed to guess the target word in three attempts. Upon reflection, I realize that I was focusing too much on the individual clues, trying to find a word that fits each description. However, I didn't consider the connections between the clues. Next time, I will try to find a word that ties together the different characteristics mentioned in the clues. I will also pay closer attention to the word's connotations and associations, as the target word seems to be related to something desirable or appealing. I will start by looking for a word that combines the ideas of "new" and "exciting", and then see how it fits with the other clues.} \\
\\ 
\texttt{Trial \#2: New plan: I failed to guess the target word in three attempts, and upon reflection, I realize that I was focusing too much on the individual clues, trying to find a word that fits each description. However, I didn't consider the connections between the clues. Next time, I will try to find a word that ties together the different characteristics mentioned in the clues. I will also pay closer attention to the word's connotations and associations, as the target word seems to be related to something desirable or appealing. } \\
\\ 
\texttt{To achieve this, I will start by looking for a word that combines the ideas of "new" and "exciting", and then see how it fits with the other clues. I will also try to find a word that is related to the idea of a "concert" or a "festival", as mentioned in the third clue. I will make sure to consider the connections between the clues and look for a word that ties them together. } \\
\\ 
\texttt{In the first trial, I will try to find a word that combines the ideas of "new" and "exciting", such as "novelty" or "spectacle". I will then see how it fits with the other clues and adjust my guess accordingly.} \\
\\ 
\texttt{Trial \#3: New plan: I failed to guess the target word in three attempts, and upon reflection, I realize that I was focusing too much on the individual clues, trying to find a word that fits each description. However, I didn't consider the connections between the clues. Next time, I will try to find a word that ties together the different characteristics mentioned in the clues. I will also pay closer attention to the word's connotations and associations, as the target word seems to be related to something desirable or appealing.} \\
\\ 
\texttt{To achieve this, I will start by looking for a word that combines the ideas of "new" and "exciting", and then see how it fits with the other clues. I will also try to find a word that is related to the idea of a "concert" or a "festival", as mentioned in the third clue. I will make sure to consider the connections between the clues and look for a word that ties them together.} \\
\\ 
\texttt{In the first trial, I will try to find a word that combines the ideas of "new" and "exciting", such as "novelty" or "spectacle". I will then see how it fits with the other clues and adjust my guess accordingly. In the second trial, I will try to find a word that is related to the idea of a "concert" or a "festival", such as "performance" or "event". I will then see how it fits with the other clues and adjust my guess accordingly.} \\
\\ 
\texttt{In the third trial, I will try to find a word that ties together the different characteristics mentioned in the clues, such as "novelty" and "performance". I will also pay closer attention to the word's connotations and associations, as the target word seems to be related to something desirable or appealing. I will make sure to consider the connections between the clues and look for a word that ties them together.} \\
\\ 
\texttt{Some specific words that I will consider in the next trial are:} \\
\\ 
\texttt{* Novelty} \\
\texttt{* Spectacle} \\
\texttt{* Performance} \\
\texttt{* Event} \\
\texttt{* Festival} \\
\texttt{* Concert} \\
\\ 
\texttt{I will also pay closer attention to the word's connotations and associations, as the target word seems to be related to something desirable or appealing.} \\
\\ 
\\ 
\texttt{You are playing a collaborative word guessing game in which you have to guess a target word that another player describes to you.} \\
\\ 
\texttt{You can make one guess at each trial. You win when you guess the target word. You lose when you cannot guess it in 3 tries.} \\
\\ 
\texttt{After each trial you will get a new hint from the other player which starts with CLUE.} \\
\\ 
\texttt{Make your guesses by just saying the word using the following form: GUESS: <a word>} \\
\\ 
\texttt{Let us start.} \\
\\ 
\\ 
\texttt{CLUE: Something that adds a unique twist to familiar things, making them more engaging and interesting.} \\
            }
        }
    }
    & \\ \\

    \theutterance \stepcounter{utterance}  
    & & & \multicolumn{4}{p{0.6\linewidth}}{
        \cellcolor[rgb]{0.9,0.9,0.9}{
            \makecell[{{p{\linewidth}}}]{
                \texttt{\tiny{[GM$\langle$P2]}}
                \texttt{GUESS: Novelty} \\
            }
        }
    }
     \\ \\

    \theutterance \stepcounter{utterance}  
    & & & \multicolumn{2}{p{0.3\linewidth}}{
        \cellcolor[rgb]{0.9,0.9,0.9}{
            \makecell[{{p{\linewidth}}}]{
                \texttt{\tiny{[GM$|$GM]}}
                \texttt{continue} \\
            }
        }
    }
    & & \\ \\

    \theutterance \stepcounter{utterance}  
    & & & \multicolumn{2}{p{0.3\linewidth}}{
        \cellcolor[rgb]{0.9,0.9,0.9}{
            \makecell[{{p{\linewidth}}}]{
                \texttt{\tiny{[GM$|$GM]}}
                \texttt{novelty} \\
            }
        }
    }
    & & \\ \\

    \theutterance \stepcounter{utterance}  
    & & & \multicolumn{2}{p{0.3\linewidth}}{
        \cellcolor[rgb]{0.9,0.9,0.9}{
            \makecell[{{p{\linewidth}}}]{
                \texttt{\tiny{[GM$|$GM]}}
                \texttt{end game} \\
            }
        }
    }
    & & \\ \\

\end{supertabular}
}

\end{document}
